\scanpage{34}
de façon précise
\[
(30)\qquad\begin{aligned}
\widetilde C_\infty/\mu_\infty&\simeq\text{Image inverse de }
\overbrace{\overline C_1/\mu_3}^{\substack{\omega\\\wr}}
\subset\mu_3^3/\mu_3=\pi_3\\
&\hspace{2cm}\text{par }\pi_\infty\longrightarrow\pi_3
\end{aligned}
\]
En vertu de la suite exacte
\diagram{\[
(31)\qquad\begin{tikzcd}[ampersand replacement=\&,column sep=large]
0\arrow[r,no head]\&\pi_\infty\arrow[r,"3"]\&\pi_\infty\arrow[r,"p_3"]\&
\pi_3\arrow[r,no head]\arrow[d,phantom,"\cup"]\&0\\
\&\&\&\omega\&
\end{tikzcd}
\]}
$\pi_\infty$ opère librement sur cette partie de $\pi_\infty$, l'image inverse de
$\omega$ par $\pi_\infty\longrightarrow\pi_3$. Ceci donne les opérations de
$\pi_\infty$ sur $E(\alpha_*\deleted{\unsure{\overline\alpha}})\simeq p_3^{-1}(\omega)$,
celles de $\Gamma$ et $\unsure{\mathfrak S_3|\mathbb G_3}$ étant induites par
leurs opérations sur $\pi_\infty\ldots$.

\uline{NB} Sur $\pi_3$ \deleted{$\unsure{\mathcal E}\times$}
$\unsure{\mathfrak S_3|\mathbb G_3}\times\Gamma$ opère via
$\unsure{\mathfrak S_3|\mathbb G_3}\times\{\pm1\}$,
\margin{opère par homothéties sur $\pi_3$}
($\Gamma\overset{\chi_3}{\longrightarrow}\{\pm1\}$).
Considérons $\unsure{\mathfrak S_3^+|\mathbb G_3^+}$ \inserted{$\simeq\mathbb Z/3\mathbb Z$}
opérant sur $\pi_3$ $(\simeq\mathbb F_3^2)$, alors on sait qu'il doit y avoir
une droite invariante -- \deleted{\unsure{et les}} en fait elle est unique -- et
$\omega$ est formé des deux éléments non nuls de cette droite\ldots
% archive foliation: 33
